% Appendix H: Cross-Architecture and Scaling

\section{Cross-Architecture and Scaling}
\label{app:cross_architecture_and_scaling}

% 本节提供 §4.2 中跨模型、跨规模分析的完整数据。

\subsection{Coder vs.\ Base Paired Analysis}
\label{app:subsec:coder_vs_base}

% Qwen2.5-Coder 与 Qwen2.5 在各任务上的 outcome 分布对比。
% FJC 归一化层位的 Pearson 相关性数据与散点图。
% 预期效果：量化 code-specialized 训练对 outcome 分布的改善及对 FJC 位置的微弱影响。

% TODO: 插入 Coder vs Base 对比数据

\subsection{Scaling Trends}
\label{app:subsec:scaling_trends}

% 0.5B → 14B 的 outcome 分布变化趋势。
% Overprocessed 非单调趋势的详细数据。
% Normalized brewing duration 在不同规模间的稳定性数据。
% 预期效果：支撑 "capability changes but scaffold does not" 的 claim。

% TODO: 插入 scaling 趋势数据与图表

\subsection{Cross-Architecture Robustness}
\label{app:subsec:cross_architecture_robustness}

% CodeLlama/Llama 和 DeepSeek-Coder 上的 brewing 存在性与 outcome fingerprint 结构。
% 预期效果：证明 brewing-to-resolution 不是 Qwen 特有的，而是 transformer 处理代码的一般属性。

% TODO: 插入跨架构实验结果
